\subsection{舒尔引理}

命题 2. —— 设 A 是环，M 与 N 是两个 A-模，f 是从 M 到 N 的一个非零同态。

\begin{enumerate}
\def\labelenumi{\alph{enumi})}
\item
  若 M 是单模，则 f 是单射。
\item
  若 \(\mathbf{N}\) 是单模，则 \(f\) 是满射。
\item
  若 M 与 N 都是单模，则 f 是同构。
\end{enumerate}

f 的核是 M 的真子模，f 的像是 N 的非零子模。

\begin{enumerate}
\def\labelenumi{\alph{enumi})}
\item
  若 M 是单模，则我们有 \(\operatorname{Ker}(f)=0\)，故 f 是单射。
\item
  若 \(N\) 是单模，则我们有 \(\operatorname{Im}(f) = N\)，故 \(f\) 是满射。
\item
  若 M 与 N 都是单模，则 f 既单又满。
\end{enumerate}

推论（舒尔引理）. —— 单模的自同态环是体。

若 M 是单 A-模，则非零环 \(\operatorname{End}_{\mathrm{A}}(\mathrm{M})\) 的每个非零元素都可逆（命题 2, c)），故 \(\operatorname{End}_{\mathrm{A}}(\mathrm{M})\) 是体。

定理 1. —— 设 K 是一个代数闭交换域，A 是一个 K-代数，M 是一个单 A-模。假设 M 作为 K 上向量空间的维数有限，或更一般地，严格小于 K 的基数。则 A-模 M 的自同态环由位似 \(\alpha_{M}\)（\(\alpha \in K\)）组成。

设 \(\mathrm{E}\) 是 A-模 \(\mathrm{M}\) 的自同态环；由命题 2 的推论它是体，且是体 \(\mathrm{K}\) 上的一个代数。若把 \(\mathrm{M}\) 看作体 \(\mathrm{E}\) 上的左向量空间，则由 II, §1, No.~13, 第 222 页的命题 25，我们有 \(\dim_{\mathrm{K}} \mathrm{M} = (\dim_{\mathrm{E}} \mathrm{M})[\mathrm{E} : \mathrm{K}]\)，从而 \(\dim_{\mathrm{K}} \mathrm{M} \geqslant [\mathrm{E} : \mathrm{K}]\) 。由于按假设 \(\dim_{\mathrm{K}} \mathrm{M} < \operatorname{Card}(\mathrm{K})\)，等式 \(\mathrm{E} = \mathrm{K} \cdot 1_{\mathrm{M}}\) 便是下面引理的推论。

引理 1. —— 设 E 是体，K 是 E 的中心的一个不等于 E 的子体。若体 K 是代数闭的，则我们有 \([E:K] \geqslant \text{Card}(K)\) 。

设 \(x\) 是 \(\mathrm{E} - \mathrm{K}\) 的一个元素，\(\mathrm{L}\) 是 \(\mathrm{E}\) 的由 \(\mathrm{K} \cup \{x\}\) 生成的（交换）子体。由于 \(\mathrm{K}\) 是代数闭的，\(x\) 在 \(\mathrm{K}\) 上是超越的。由 VII, §2, No.~3, 第 10 页定理 2 及第 11 页，元素 \((x - \alpha)^{-1}\)（属于 \(\mathrm{L}\)，其中 \(\alpha\) 取遍 \(\mathrm{K}\)）在 \(\mathrm{K}\) 上线性无关。因此我们有 \([\mathrm{E} : \mathrm{K}] \geqslant [\mathrm{L} : \mathrm{K}] \geqslant \operatorname{Card}(\mathrm{K})\) 。

例. —— *设 A 是一个由可数族元素生成的 C-代数；它在 C 上有可数维。设 M 是一个单 A-模；

它是单生成的，故在 \(\mathbf{C}\) 上具有可数基。由于体 \(\mathbf{C}\) 不是可数的（Gen.~Top., IV, §4, No.~1, 第 44 页），我们有 \([\mathrm{M}:\mathbf{C}] < \mathrm{Card}(\mathbf{C})\) 。因此 A-模 M 的自同态就是位似 \(\alpha_{\mathrm{M}}\)（\(\alpha \in \mathbf{C}\)）。这特别适用于 A 是 \(\mathbf{C}\) 上有限维李代数的泛包络代数的情形（Lie, I, §2, No.~7, 第 21 页，推论 3）。*

推论 1. —— 保留定理 1 的假设，并再设代数 A 是交换的。则 M 在 K 上的维数为 1。

由于环 A 交换，\(a_{M}\) 对每个 \(a \in A\) 都是 A-模 M 的自同态。于是由定理 1，我们有 \(A_{M} = K \cdot 1_{M}\)，且 M 是单 K-模，即体 K 上的 1 维向量空间。

推论 2. —— 设 K 是交换域，A 是一个 K-代数，M 是一个 A-模。假设对 K 的每个扩张 L，\(\mathrm{A}_{(\mathrm{L})}\) -模 \(\mathrm{M}_{(\mathrm{L})}\) 都是单模。则 M 的自同态环由位似 \(\alpha_{M}\)（\(\alpha \in K\)）组成。

设 I 是一个基数严格大于 M 在 K 上的维数的集合（例如 M 的子集组成的集合）。设 L 是体 \(\mathrm{K}((\mathrm{X}_{i})_{i\in\mathrm{I}})\) 的一个代数闭包（V, §4, No.~3, 第 23 页，定理 2）。取 L 上的一个 K-线性型 \(\varphi\)，使 \(\varphi(1)=1\)，并设 \(v:M_{(\mathrm{L})}\to M\) 是由 \(v(\alpha\otimes m)=\varphi(\alpha)m\) 刻画的 K-线性映射。设 u 是 M 的一个自同态。\(M_{(\mathrm{L})}\) 在 L 上的维数等于 M 在 K 上的维数，且严格小于 L 的基数。由定理 1，\(1_{L}\otimes u\)——\(A_{(\mathrm{L})}\) -模 \(M_{(\mathrm{L})}\) 的自同态——形如 \(\lambda\otimes1_{M}\)（\(\lambda\in L\)）。对每个 \(x\in M\)，我们有

\[
\begin{array}{c} u (x) = v \bigl (1 \otimes u (x) \bigr) = v \bigl ((1 _ {\mathrm{L}} \otimes u) (1 \otimes x) \bigr) \\ = v \bigl ((\lambda \otimes 1 _ {\mathrm{M}}) (1 \otimes x) \bigr) = v (\lambda \otimes x) = \varphi (\lambda) x, \end{array}
\]

于是 u 是位似 \(\varphi(\lambda)_{\mathrm{M}}\) 。

